%! Logarithmic Functions — Unit 2, Lesson 2.11 — Guided Notes (Answer Key)
\documentclass[10pt]{article}
\usepackage{precalculus-article}
\usepackage{precalculus-key}   % pulls in -boxes; do NOT also load -boxes
\newcommand{\TallMath}[1]{$\displaystyle #1\rule[-1.4em]{0pt}{3.2em}$}

\begin{document}

\pageheader{Unit 2, Lesson 2.11}{Guided Notes --- Answer Key}
\namedateperiod

\vspace{0.08in}

\begin{objectivebox}
\textbf{LO 2.11.A:} Identify key characteristics of logarithmic functions.

By the end of this lesson I will be able to:
\begin{itemize}[leftmargin=1.5em,itemsep=3pt,topsep=3pt]
  \item State the domain, range, and vertical asymptote of $f(x) = a\log_b x$.
  \item Describe the end behavior, monotonicity, and concavity of a logarithmic function.
  \item Use the additive transformation identification test to decide if a function is logarithmic.
\end{itemize}
\end{objectivebox}

\vspace{0.06in}

\begin{vocabbox}
\termblanklong{logarithmic function (general form)}
\ansline{$f(x) = a\log_b x$, where $b > 0$, $b \ne 1$, $a \ne 0$; the inverse of $g(x) = b^x$.}

\termblanklong{domain of a logarithmic function}
\ansline{All real numbers greater than zero: $(0, \infty)$.}

\termblanklong{vertical asymptote}
\ansline{The line $x = 0$; the graph of $a\log_b x$ approaches but never touches it.}

\termblanklong{end behavior}
\ansline{How $f(x)$ behaves as $x \to 0^+$ and as $x \to \infty$; determined by $a$ and $b$.}

\termblanklong{monotone increasing / decreasing}
\ansline{Always increasing or always decreasing on its full domain; no local extrema.}

\termblanklong{concave up / concave down}
\ansline{The graph curves in one direction only; logarithms have no inflection point.}

\termblanklong{additive transformation identification test}
\ansline{If $g(x) = f(x+k)$ has inputs proportional over equal output intervals, then $f$ is logarithmic.}
\end{vocabbox}

\vspace{0.06in}

\begin{hookbox}
\textbf{Richter Scale Context:} The Richter scale is defined by $M = \log_{10}(I)$, where
$I$ is earthquake intensity.

\vspace{4pt}
If intensity $I = 10$, what is the magnitude? \ans{1}
If intensity $I = 100$, magnitude $= $ \ans{2}
If intensity $I = 1000$, magnitude $= $ \ans{3}

\vspace{4pt}
Each time $I$ multiplies by \ans{10}, the magnitude $M$ increases by \ans{1}.

\vspace{4pt}
What inputs $I$ are physically meaningful for this function? \ansline{$I > 0$; intensity must be a positive number.}
\end{hookbox}

\vspace{0.06in}

\begin{notesbox}{Domain and Range (EK 2.11.A.1)}
For $f(x) = a\log_b x$ in general form:

\vspace{6pt}
\begin{center}
\renewcommand{\arraystretch}{1.8}
\begin{tabular}{|l|l|}
\hline
\rowcolor{sky}
\textbf{Characteristic} & \textbf{Value} \\
\hline
Domain & \ans{$x > 0$, i.e., $(0, \infty)$} \\
\hline
Range  & \ans{All real numbers, i.e., $(-\infty, \infty)$} \\
\hline
Vertical asymptote & $x = $ \ans{$0$} \\
\hline
\end{tabular}
\end{center}

\vspace{6pt}
Why can't we evaluate $\log_b 0$ or $\log_b(-3)$? \ansline{There is no exponent $t$ such that $b^t = 0$ or $b^t = -3$; logarithms are undefined for non-positive inputs.}
\end{notesbox}

\vspace{0.06in}

\begin{notesbox}{Characteristics from the Inverse (EK 2.11.A.2)}
Because $f(x) = \log_b x$ is the inverse of $g(x) = b^x$, it inherits structural properties.
Complete the comparison table.

\vspace{6pt}
\begin{center}
\renewcommand{\arraystretch}{1.8}
\begin{tabular}{|l|l|l|}
\hline
\rowcolor{sky}
\textbf{Property} & \textbf{Exponential $g(x) = b^x$} & \textbf{Logarithmic $f(x) = \log_b x$} \\
\hline
Always increasing ($b>1$)? & Yes & \ans{Yes} \\
\hline
Always decreasing ($0<b<1$)? & Yes & \ans{Yes} \\
\hline
Has local extrema (natural domain)? & No & \ans{No} \\
\hline
Has an inflection point? & No & \ans{No} \\
\hline
\end{tabular}
\end{center}

\vspace{6pt}
So logarithmic functions are always \ans{increasing} or always \ans{decreasing}, and their
graphs are either always \ans{concave up} or always \ans{concave down}.
\end{notesbox}

\vspace{0.06in}

\begin{notesbox}{End Behavior (EK 2.11.A.4)}
Fill in the end behavior for each case. The sign of $a$ and the value of $b$ both matter.

\vspace{6pt}
\begin{center}
\renewcommand{\arraystretch}{2.0}
\begin{tabular}{|l|l|l|}
\hline
\rowcolor{sky}
\textbf{Conditions} & \TallMath{\lim_{x \to 0^+} a\log_b x} & \TallMath{\lim_{x \to \infty} a\log_b x} \\
\hline
$b > 1$, $a > 0$ & \ans{$-\infty$} & \ans{$+\infty$} \\
\hline
$b > 1$, $a < 0$ & \ans{$+\infty$} & \ans{$-\infty$} \\
\hline
$0 < b < 1$, $a > 0$ & \ans{$+\infty$} & \ans{$-\infty$} \\
\hline
\end{tabular}
\end{center}
\end{notesbox}

\vspace{0.06in}

\begin{notesbox}{Additive Transformation Identification Test (EK 2.11.A.3)}
\textbf{Key Idea:} If $g(x) = f(x + k)$ for $k \ne 0$, then $g$ is an additive
transformation of $f$.

\begin{itemize}[leftmargin=1.5em, itemsep=4pt, topsep=4pt]
  \item If the \emph{inputs} of $g(x) = f(x+k)$ are proportional over equal-length output
        intervals, then $f$ is logarithmic.
  \item If the inputs of $f$ itself are proportional over equal-length output intervals,
        that confirms $f$ is logarithmic.
\end{itemize}

\vspace{6pt}
\textit{Example.} Below is a table for $f(x) = \log_3 x$. Check: as the output increases
by 1, does the input multiply by a constant?

\vspace{4pt}
\begin{center}
\renewcommand{\arraystretch}{1.8}
\begin{tabular}{|c|c|c|}
\hline
\rowcolor{sky}
$x$ (input) & $f(x)$ (output) & Ratio of consecutive inputs \\
\hline
$1$  & $0$ & --- \\
\hline
$3$  & $1$ & $3/1 = $ \ans{3} \\
\hline
$9$  & $2$ & $9/3 = $ \ans{3} \\
\hline
$27$ & $3$ & $27/9 = $ \ans{3} \\
\hline
$81$ & $4$ & $81/27 = $ \ans{3} \\
\hline
\end{tabular}
\end{center}

\vspace{6pt}
The inputs are proportional (each multiplied by \ans{3}), so $f$ \ans{is}
(is / is not) logarithmic.
\end{notesbox}

\end{document}
